\section{Conclusion}
\label{sec:conclusion}

This paper asks whether transparent graph diagnostics can decide when node classification should use a tuned graph-model family rather than a tuned feature-only MLP. In a frozen prospective benchmark spanning 11 datasets, ten seeds, 770 model records, and 110 diagnostic records, the historical combined rule reaches 55.5\% selection accuracy and 7.46 percentage points of full-set regret. Always-graph reaches 80.9\% and 0.26 points, while validation selection reaches 91.8\% and 0.22 points. No predeclared comparison establishes an advantage for the combined rule after Holm correction. Under this protocol, the evaluated simple diagnostics therefore provide no stable incremental decision value.

The negative decision result is not evidence that graph structure is irrelevant. Degree-preserving randomization lowers GCN accuracy by 45.5, 29.9, and 17.3 points on Cora, CiteSeer, and PubMed, respectively. Because the intervention changes several forms of edge organization and covers only three dataset clusters, it does not isolate homophily or support a universal causal claim. It instead supplies the paper's main bounded insight: graph structure can carry substantial predictive value even when low-dimensional diagnostics fail to characterize that value reliably.

Future work should evaluate richer or learned representations of graph utility under the same prospective safeguards, including matched trivial baselines, explicit abstention, regret, dataset-level inference, and immutable provenance. The scoped fixed-degree analysis offers mechanistic context for such work, but the empirical evidence argues against treating a single structural statistic as a general graph-vs-MLP selection rule.
